\FloatBarrier
\section{Torus coil sharing with openings retained}
\label{sec:coils}

Upcoming in GEMC2 6.2 and CLAS12 GEMC Systems 0.4, the six copper torus coils share one unchanged engineering
mesh. coil1 is renamed coil, while coil2 through coil6 keep their sector suffixes and become five passive
copies. All openings, rounded corners, thickness, G4\_Cu material, and colors are retained. Both source and
runtime collections are synchronized in the two owning repositories; the six engineering originals are
archived separately.

Every original has 508 vertices and 1,016 triangles in one closed manifold component, with genus one, no
boundary edges, and no detected self-intersections. Exact 60-degree sector rotations give a one-to-one vertex
correspondence with maximum residual 0.000224728 mm, consistent with binary STL coordinate precision. Four
exports have identical mapped facet connectivity. coil2 and coil3 each use the opposite diagonal on one planar
side quadrilateral; the builder verifies the same four corners, perimeter, and plane. This changes the
triangulation rather than the represented surface.

The canonical coil STL is byte-identical to coil1. All 1,016 facets, including every surface around the large
central opening, remain. The source is already a clean tessellated ring; sharing achieves the reduction
without further contour simplification or hole filling. The copy placements use integer sector angles and
millimetre positions with three decimal places. No empirical translation or contact relief is introduced.

Unique stored and built coil facets fall from 6,096 to 1,016, a reduction of 5,080 (83.33\%). Active STL
storage falls from 305,304 to 50,884 bytes. All six physical coils remain, with 6,096 placed facets. Total
signed copper volume changes from 691,360,208.654 to 691,360,128.713 cubic millimetres (-0.000011563\%)
because the canonical coordinates replace separately rounded exports. The complete magnets collection now has
28 STL sources and 69,750 unique facets for 39 CAD placements.

GEMC3 keeps the sources and copies together in magnets\_cad. GEMC2 uses 28 CAD sources and 11 native
magnets\_copies rows per inherited run, including the six existing jacket copies. The supplied cards already
load this companion bank. All nine GEMC3, GEMC2 SQLite/CAD, and GEMC2 TEXT/CAD construction checks pass across
runs 11, 3029, and 4763. Semantic comparison matches all three native solenoid configurations and the three
39-row CAD banks. Unrelated native geometry, materials, and DC banks are unchanged. The builder reproduces the
canonical mesh, health report, and all three placement catalogs byte for byte.

The full run-11 assembly contains native DC, both DC CAD groups, solenoid, magnets, and native/CAD beamline:
139 volumes including the world. Matched baseline and final fixtures use seed 12345 and 100,000 Geant4 surface
samples per tested volume. Readout is disabled in these geometry-only fixtures. Existing orientation and
overlap warnings remain; the baseline reports 22 ordered overlap pairs and the final fixture reports 22. No
new sampled overlap pair is introduced by sharing the coil. Sampling is not a mathematical clearance proof.

The timing fixture isolates the six coil placements at run 11. The baseline builds six independent solids; the
optimized fixture builds one solid and reuses it through five copies. One fixed GEMC3 dev executable reports
geant4-11-03-patch-02 [MT]   (25-April-2025). It uses the serial run manager, FTFP\_BERT, default production
cuts, one execution thread, no field, no optical physics, no GUI, no digitization, and no event output.
Five-GeV/c electrons, protons, positive pions, and neutrons start at the origin, uniformly in solid angle from
5 to 70 degrees with full azimuth. There is no target, luminosity normalization, or pileup. Five
alternating-order repetitions use 2,000 matched LUND primaries per particle and seeds 12345--12349. Fifty
timed processes and two isolated overlap checks are scheduled. Every successful tracking process reports 2,000
events; 7 serial processes fail as recorded below. Statistics use only seeds for which both fixtures complete;
a failed pair is excluded from both timing samples. Sleep prevention remains enabled throughout.

An initial tasking-mode attempt stopped after twenty processes when the optimized 2,000-event electron run at
seed 12349 raised fatal PART122: Cu62 was already registered in the Geant4 particle table during a
photonuclear interaction. This is a particle-table error, rather than a reported geometry overlap. Both
fixtures complete those electron primaries in a serial-mode probe, but the full serial repeat also fails on
the original proton geometry at seed 12349 with duplicate Co60[58.590] registration. Thus serial execution
does not resolve the issue, and it affects the reference geometry as well. The scheduled serial processes are
continued after failures, with the same executable hash. All failures, raw logs, and probe inputs are
archived. Tasking and serial results are not pooled. These failures prevent an unconditional
tracking-performance conclusion.

Median zero-event process wall time falls from 0.18605 to 0.13280 s (1.401 times faster). Median process CPU
time changes from 0.100 to 0.090 s, and peak RSS from 63.17 to 59.25 MiB. Tracking entries below are median
internal transport wall times; before/after ratios above one indicate faster tracking. The paired range
retains seed dependence. These isolated coil measurements do not establish a full-CLAS12 speedup or
detector-response equivalence, and the tiny geometric differences can alter stochastic trajectories. Paired
startup ratios range from 1.004 to 1.451. Ranges that span one do not establish a consistent gain in this
sample. The interrupted tasking attempt completed its ten zero-event measurements: median loading wall times
were 0.68288 and 0.67769 s, essentially unchanged, with peak RSS 121.70 and 117.52 MiB. Performance depends on
execution mode; the table below uses serial mode throughout.

\begin{table}[htbp]
\centering
\setlength{\tabcolsep}{6pt}
\caption{Torus coil tracking; completed pairs of 2,000-event processes only.}
\label{tab:coils}
\begin{tabular}{lrrrrr}
\toprule
Particle & Pairs & Before (s) & After (s) & Ratio & Paired ratio range \\
\midrule
Electron & 5 & 6.41057 & 5.80693 & 1.104 & 1.043--1.168 \\
Proton & 4 & 4.97728 & 3.90553 & 1.274 & 1.207--1.451 \\
Positive pion & 2 & 4.99705 & 3.99851 & 1.250 & 1.211--1.289 \\
Neutron & 3 & 3.03117 & 2.02905 & 1.494 & 1.303--1.562 \\
\bottomrule
\end{tabular}
\end{table}

\begin{figure}[htbp]
\centering
\includegraphics[width=\linewidth]{figures/coil_comparison.pdf}
\caption{
All six engineering coil exports aligned to sector 1, with their full ring profiles and central opening
retained. The right panel shows maximum matched-vertex residuals under exact nominal sector rotations; the
largest is 0.225 micrometres. The canonical mesh is unchanged, so an additional before/after surface overlay
would coincide.}
\label{fig:coils}
\end{figure}

data/magnets\_coils and magnets\_coils\_archive.zip retain engineering mesh analysis, the shared-mesh proof,
material-volume changes, database comparisons, all nine construction logs, both dense assembly fixtures,
repeated timing inputs, matched primaries, raw logs, and executable hash. The archive includes the six
originals, final collected geometry and meshes, catalogs, builder, benchmark, and validation code. The
executable is not redistributed. Earlier study archives remain unchanged.
